% Generated draft: requires mathematical and source review.
\documentclass[11pt]{article}
\usepackage[margin=25mm]{geometry}
\usepackage{fontspec}
\setmainfont{Latin Modern Roman}
\usepackage{amsmath,amssymb,mathtools,unicode-math}
\setmathfont{Latin Modern Math}
\usepackage{enumitem,needspace,xurl,hyperref}
\setlength{\parindent}{0pt}
\setlength{\parskip}{5pt}
\setlength{\emergencystretch}{3em}
\newcommand{\N}{\mathbb{N}}
\newcommand{\Z}{\mathbb{Z}}
\newcommand{\Q}{\mathbb{Q}}
\newcommand{\R}{\mathbb{R}}
\newcommand{\C}{\mathbb{C}}
\newcommand{\F}{\mathbb{F}}
\newcommand{\A}{\mathbb{A}}
\newcommand{\PP}{\mathbb{P}}
\newcommand{\E}{\mathbb{E}}
\newcommand{\eps}{\varepsilon}
\newcommand{\dd}{\,\mathrm{d}}
\begin{document}
\section*{UnsolvedMath ID: 10900051}
\textbf{Source problem code:} catalogue entry
\subsection{Definitions and mathematical statement}
A number field \(k\) is totally real if every embedding \(k\hookrightarrow\C\) has image in \(\R\). For an oriented complete hyperbolic \(3\)-manifold \(M=\mathbb H^3/\Gamma\), let \(\widetilde\Gamma\) be the inverse image of its holonomy group in \(\operatorname{SL}(2,\C)\). Its ordinary trace field is
\[
 t(M)=\Q\bigl(\operatorname{tr}\gamma:\gamma\in\widetilde\Gamma\bigr).
\]
The invariant trace field is the generally different field
\[
 k(M)=\Q\bigl(\operatorname{tr}(\gamma^2):\gamma\in\widetilde\Gamma\bigr).
\]

Manning, relaying an old Neumann--Reid question, asks whether for every number field \(k\) that is not totally real there is a hyperbolic \(3\)-manifold with \emph{trace field} \(k\). With ordinary trace-field terminology, the target is \(t(M)\cong k\). The source statement does not explicitly impose finite volume or say whether equality is of abstract fields or of specified subfields of \(\C\).

The supplied summary instead identifies the question with Neumann's conjecture for the \emph{invariant} trace field of each non-real concrete number field. A concrete field includes a specified embedding in \(\C\), so this is not identical wording. Ordinary versus invariant trace field, abstract versus concrete realization, and finite-volume hypotheses must therefore be resolved before selecting a single formulation.

The summary reports realization of non-real multiquadratic fields, experimental support and an unrestricted arithmetic-orbifold realization assertion for non-totally-real fields. These are retained as supplied literature claims. In particular, the last assertion needs checking: arithmetic Kleinian-lattice trace fields have restrictive complex-place hypotheses, so its unrestricted scope is not adopted as a theorem. Orbifold realization and manifold realization must also be distinguished.
\subsection{Short English statement}
Can every number field that is not totally real occur as a hyperbolic \(3\)-manifold trace field? The source says trace field, while its summary switches to invariant trace fields and concrete embeddings.
\subsection{Sources}
\begin{itemize}
\item Manning, attributed item 7.3, described as an old Neumann--Reid question, in Delp, Hoffoss and Manning, \emph{Problems In Groups, Geometry, and Three-Manifolds} (2015), supplied PDF page 6. \url{https://arxiv.org/pdf/1512.04620#page=6}; original-status link \url{https://arxiv.org/abs/1512.04620}. Recorded access 2026-07-29; extraction source-tex checked against rendered PDF text, as reported in the record.
\item The supplied summary cites \emph{An experimental investigation of Neumann's conjecture}, \url{https://arxiv.org/abs/1609.08719}, mentions discussion with Reid, Neumann's realizing-arithmetic-invariants lectures (Columbia volume/conference, 2006--2009), and a Reid--Neumann multiquadratic-field observation. Complete lecture and partial-result citations are absent; no author list is reconstructed for the arXiv item.
\item Local UnsolvedMath registry, \texttt{data/unsolvedmath/problems.json}, permanent ID \(10900051\), exact source code \texttt{catalogue entry}; supplied partially\_solved/PARTIAL-PROGRESS status.
\end{itemize}
\end{document}
